\documentclass{article}
\usepackage{amsmath, amssymb, amsthm}
\usepackage{hyperref}
\usepackage{geometry}
\geometry{margin=1in}

\title{Erd\H{o}s Problem \#1104}
\date{Last updated: 2025-10-26}
\author{Source: erdosproblems.com}

\begin{document}
\maketitle

\noindent\textbf{Status:} OPEN \quad \textbf{No prize}

\noindent\textbf{Tags:} graph theory, chromatic number

\medskip
\noindent\textbf{Problem Statement:}

Let $f(n)$ be the maximum possible chromatic number of a triangle-free graph on $n$ vertices. Estimate $f(n)$.

\bigskip
\noindent\textbf{Remarks:}

The best bounds available are\[(1-o(1))(n/\log n)^{1/2}\leq f(n) \leq (2+o(1))(n/\log n)^{1/2}.\]The upper bound is due to Davies and Illingworth \cite{DaIl22}, the lower bound follows from a construction of Hefty, Horn, King, and Pfender \cite{HHKP25}.

One can ask a similar question for the maximum possible chromatic number of a triangle-free graph on $m$ edges. Let this be $g(m)$. Davies and Illingworth \cite{DaIl22} prove\[g(m) \leq (3^{5/3}+o(1))\left(\frac{m}{(\log m)^2}\right)^{1/3}.\]Kim \cite{Ki95} gave a construction which implies $g(m) \gg (m/(\log m)^2)^{1/3}$.

The function $f(n)$ is the inverse to the function $h_3(k)$ considered in [1013].

A generalisation of $f(n)$ is considered in [920].

\bigskip
\noindent\textbf{References:}

[DaIl22] Davies, Ewan and Illingworth, Freddie, The {$\chi$}-{R}amsey problem for triangle-free graphs. SIAM J. Discrete Math. (2022), 1124--1134.
 
 [HHKP25] Z. Hefty, P. Horn, D. King, and F. Pfender, Improving $R(3,k)$ in just two bites. arXiv:2510.19718 (2025).
 
 [Ki95] Kim, J. H., The Ramsey number $R(3,t)$ has order of magnitude $t^2/\log t$. Random Structures and Algorithms (1995), 173-207.

\bigskip
\noindent\small{Source: \url{https://www.erdosproblems.com/1104}}
\end{document}
